\begin{definition}\label{sga3:XXI:3.6.2}
A subset \(C\subset V\) described by the construction of
Lemma~\ref{sga3:XXI:3.6.1} is called, here, a \emph{chamber} of \(V\).
\end{definition}

\begin{definition}\label{sga3:XXI:3.6.3}
A hyperplane \(H\) of \(V\) is called a \emph{wall} of \(C\) if
\[
 H\cap(\overline C\setminus C)
\]
contains a nonempty open subset of \(H\).
\end{definition}

\begin{remark}\label{sga3:XXI:3.6.4}
For a convex subset, the closure can be described without reference
to the topology of \(V\): it is the set of endpoints of all open
segments contained in the given subset.
\end{remark}

\begin{lemma}\label{sga3:XXI:3.6.5}
Under the conditions of Lemma~\ref{sga3:XXI:3.6.1},
\[
 \overline C=\{x\in V\mid f_i(x)\ge0\}.
\]
The walls of \(C\) are the hyperplanes \(f_i^{-1}(0)\).
\end{lemma}

\begin{proof}
The first assertion is clear. The second follows from
\[
 \overline C\setminus C\subset\bigcup_i f_i^{-1}(0)
\]
and from the fact that the \(f_i^{-1}(0)\) are visibly walls of
\(C\).
\end{proof}

\begin{proposition}\label{sga3:XXI:3.6.6}
Let \(C\) be a chamber of \(V\). If
\[
 H_i,\qquad i=1,2,\ldots,n,
\]
are the distinct walls of \(C\), then, for every system of linear
forms \(\{u_i\}\) such that \(H_i=u_i^{-1}(0)\), there exist
\(\varepsilon_i\in\{-1,+1\}\) for which
\[
 C=\{x\in V\mid\varepsilon_i u_i(x)>0\}.
\]
For every wall \(H\) of \(C\),
\[
 H\cap C=\varnothing,
 \qquad
 H\cap\overline C=\overline{C_H},
\]
where \(C_H\) is a chamber in \(H\). The walls of \(C_{H_i}\) are
the \(H_j\cap H_i\), for \(j\ne i\).
\end{proposition}

\begin{proof}
This follows immediately from the lemma.
\end{proof}

\begin{definition}\label{sga3:XXI:3.6.7}
The \(C_{H_i}\) are called the \emph{faces} of \(C\).
\end{definition}

Let now
\[
 \cR=(M,M^\ast,R,R^\ast)
\]
be a root datum, and put
\[
 V_{\mathbb{R}}^\ast=M^\ast\otimes\mathbb{R}.
\]

\begin{definition}\label{sga3:XXI:3.6.8}
\footnote{Editors' note: The definition of the hyperplanes
\(\mathcal{H}_\alpha\) was added here.}
For every \(\alpha\in R\), put
\[
 \mathcal{H}_\alpha
   =\{x\in V_{\mathbb{R}}^\ast\mid(x,\alpha)=0\}.
\]
Write
\[
 X=V_{\mathbb{R}}^\ast
   \setminus\bigcup_{\alpha\in R}\mathcal{H}_\alpha.
\]
For every \(x\in X\), put
\[
 R_+(x)=\{\alpha\in R\mid(x,\alpha)>0\}.
\]
For every system of positive roots \(R_+\), put
\[
 \mathcal{C}(R_+)
   =\{x\in V_{\mathbb{R}}^\ast
       \mid (x,\alpha)>0\text{ for every }\alpha\in R_+\}.
\]
\end{definition}
